\section{Modul 10. Logika dan geser bit}

\subsection{Tujuan}

Setelah menyelesaikan modul ini, mahasiswa mampu:
\begin{enumerate}[nosep]
    \item menjelaskan operasi \texttt{AND}, \texttt{OR},
          \texttt{XOR}, dan \texttt{NOT} pada tingkat bit;
    \item memakai mask untuk menyalakan, memadamkan, atau membalik
          bit tertentu;
    \item memakai \texttt{TEST} untuk memeriksa satu bit tanpa
          mengubah operand;
    \item menjelaskan \texttt{SHL} dan \texttt{SHR} serta peran
          \texttt{CF};
    \item menulis program yang memakai mask dan pergeseran bit.
\end{enumerate}

\subsection{Dasar singkat}

\subsubsection*{Operasi logika}

Operasi logika memperlakukan operand sebagai kumpulan bit yang
berdiri sendiri. Hasil setiap posisi hanya bergantung pada pasangan
bit pada posisi itu. Tabel kebenaran \texttt{AND}, \texttt{OR}, dan
\texttt{XOR} ditunjukkan pada Tabel~\ref{tab:m10-kebenaran};
\texttt{NOT} membalik setiap bit.

\begin{table}[H]
    \centering
    \caption{Tabel kebenaran operasi logika.}
    \label{tab:m10-kebenaran}
    \small
    \begin{tabular}{@{}cccc@{}}
        \toprule
        \textbf{Bit A} & \textbf{Bit B} & \textbf{AND} &
        \textbf{OR} / \textbf{XOR} \\
        \midrule
        0 & 0 & 0 & 0 / 0 \\
        0 & 1 & 0 & 1 / 1 \\
        1 & 0 & 0 & 1 / 1 \\
        1 & 1 & 1 & 1 / 0 \\
        \bottomrule
    \end{tabular}
\end{table}

\texttt{AND} menghasilkan 1 hanya bila kedua bit bernilai 1.
\texttt{OR} menghasilkan 1 bila salah satu atau kedua bit bernilai 1.
\texttt{XOR} menghasilkan 1 bila kedua bit berbeda. \texttt{NOT}
membalik seluruh bit operand dan tidak mengubah flag. Contoh dengan
\texttt{AL=0A5h}:

\begin{lstlisting}[numbers=none]
AND AL, 0F0h        ; AL = A0h, nibble bawah dipadamkan
OR  AL, 0Fh         ; AL = AFh, nibble bawah dinyalakan
XOR AL, 0FFh        ; AL = 50h, seluruh bit dibalik
NOT AL              ; AL = AFh, seluruh bit dibalik lagi
\end{lstlisting}

\subsubsection*{Mask dan TEST}

Mask adalah pola bit yang dirancang untuk menyentuh posisi tertentu
saja. Bit mask bernilai 1 mempertahankan atau menyalakan bit,
sedangkan bit mask bernilai 0 memadamkannya. Pola pemakaian mask
dirangkum pada Tabel~\ref{tab:m10-mask}.

\begin{table}[H]
    \centering
    \caption{Pemakaian mask pada satu bit tertentu.}
    \label{tab:m10-mask}
    \small
    \begin{tabular}{@{}lll@{}}
        \toprule
        \textbf{Kehendak} & \textbf{Instruksi} & \textbf{Mask} \\
        \midrule
        menyalakan bit & \texttt{OR} & 1 pada posisi bit \\
        memadamkan bit & \texttt{AND} & 0 pada posisi bit \\
        membalik bit & \texttt{XOR} & 1 pada posisi bit \\
        memeriksa bit & \texttt{TEST} & 1 pada posisi bit \\
        \bottomrule
    \end{tabular}
\end{table}

\texttt{TEST} melakukan operasi \texttt{AND} tanpa menyimpan hasil.
Kedua operand tidak berubah; yang diperbarui hanya flag. Bila hasil
\texttt{AND} bernilai nol, \texttt{ZF} bernilai 1. Karena itu,
\texttt{TEST AL,80h} diikuti \texttt{JZ} melompat bila bit 7 padam,
dan tidak melompat bila bit 7 menyala.

\subsubsection*{Pergeseran SHL dan SHR}

\texttt{SHL} menggeser seluruh bit ke kiri dan mengisi bit paling
kanan dengan nol. \texttt{SHR} menggeser seluruh bit ke kanan dan
mengisi bit paling kiri dengan nol. Pada keduanya, bit yang terdesak
keluar ditampung pada \texttt{CF}. Pergeseran satu bit ke kiri setara
dengan mengalikan dua, sedangkan pergeseran satu bit ke kanan setara
dengan membagi dua untuk bilangan tidak bertanda. Bila jumlah
pergeseran lebih dari satu, jumlah itu disimpan pada \texttt{CL}.

Contoh: \texttt{AL=09h} digeser satu bit ke kanan menghasilkan
\texttt{AL=04h} dan \texttt{CF=1}, karena bit 0 yang bernilai 1
keluar. Hasil \texttt{04h} masih satu digit sehingga dapat dicetak
setelah ditambah \texttt{30h}. Hasil yang lebih dari 9 tidak dapat
dicetak dengan cara itu dan harus diperiksa pada register.

Rujukan singkat: Gunawan, Bagian 7; S'to, Bab X bagian logika dan
pergeseran; praktikum \texttt{bab4.tex} bagian logika dan
pergeseran; buku ajar, \texttt{chapters/11}.

\subsection{Alat dan bahan}

\begin{itemize}[nosep]
    \item Komputer dengan emu8086 terpasang.
    \item DOSBox dengan TASM, TLINK, dan DEBUG.
    \item Berkas \texttt{code/m10-mask.asm} dan
          \texttt{code/m10-geser.asm}.
    \item Lembar jurnal praktikum.
\end{itemize}

\subsection{Contoh program}

Program pertama menyalakan bit 7 dengan \texttt{OR}, mempertahankan
nibble atas dengan \texttt{AND}, lalu memeriksa bit 7 dengan
\texttt{TEST}. Kalimat ``ya'' tercetak bila bit 7 menyala; bila
padam, yang tercetak ``tidak''.

\lstinputlisting[
    caption={Mask dengan \texttt{OR} dan \texttt{AND}, lalu
    \texttt{TEST} satu bit.},
    label={lst:m10-mask}
]{code/m10-mask.asm}

Program kedua menggeser \texttt{09h} satu bit ke kanan dengan
\texttt{SHR}. Hasil \texttt{04h} dapat dilihat pada \texttt{AL},
lalu dicetak sebagai karakter \texttt{'4'}.

\lstinputlisting[
    caption={Pergeseran satu byte dengan \texttt{SHR}.},
    label={lst:m10-geser}
]{code/m10-geser.asm}

\subsection{Langkah kerja}

\begin{enumerate}
    \item Siapkan folder kerja, lalu salin kedua berkas \texttt{.asm}
          ke dalamnya.
    \item Rakit dan tautkan program pertama. Jalankan kedua perintah
          berikut pada baris terpisah:
          \begin{quote}
          \texttt{tasm m10-mask.asm}\\
          \texttt{tlink m10-mask.obj}
          \end{quote}
          Jalankan hasilnya dan catat kalimat yang tercetak.
    \item Ubah mask pada \texttt{OR} dari \texttt{10000000b} menjadi
          \texttt{00000000b}. Rakit ulang, jalankan, dan jelaskan
          mengapa keluarannya berubah.
    \item Jalankan program pertama dengan \emph{single step}. Catat
          \texttt{AL} sesudah \texttt{OR}, sesudah \texttt{AND}, dan
          nilai \texttt{ZF} sesudah \texttt{TEST}.
    \item Rakit dan tautkan program kedua dengan perintah yang sama,
          lalu jalankan. Catat karakter yang tercetak.
    \item Jalankan program kedua dengan \emph{single step}. Catat
          \texttt{AL} dan \texttt{CF} sesudah \texttt{SHR}. Ubah
          nilai awal menjadi \texttt{0Ch}, rakit ulang, dan pastikan
          keluarannya berubah.
\end{enumerate}

\subsection{Tugas}

\begin{enumerate}
    \item Jelaskan perbedaan hasil \texttt{AND}, \texttt{OR}, dan
          \texttt{XOR} untuk pasangan bit 0 dan 1.
    \item Jelaskan cara memadamkan bit 3 dengan \texttt{AND} dan
          menyalakannya dengan \texttt{OR}. Tentukan mask yang
          dipakai.
    \item Jelaskan mengapa \texttt{TEST} tidak mengubah operand dan
          bagaimana \texttt{ZF} dipakai sesudahnya.
    \item Jelaskan kerja \texttt{SHL} dan \texttt{SHR}, termasuk peran
          \texttt{CF}. Tentukan hasil \texttt{SHR AL,1} bila
          \texttt{AL=0A5h}.
    \item Jelaskan mengapa \texttt{SHR} hanya tepat dipakai sebagai
          pembagian dua untuk bilangan tidak bertanda.
\end{enumerate}

\subsection{Hasil yang dikumpulkan}

\begin{enumerate}
    \item Kedua berkas \texttt{.asm} beserta versi modifikasinya.
    \item Tangkapan layar keluaran ``ya''/``tidak'' dan keluaran
          program kedua.
    \item Tabel jejak berisi \texttt{AL}, \texttt{ZF}, dan
          \texttt{CF} sesudah setiap operasi.
    \item Lembar jurnal berisi jawaban langkah kerja dan tugas.
\end{enumerate}
